\section{Secure Boot - Later stages}

\subsection{U-Boot components}

\begin{frame}{Early boot}
  \begin{itemize}
    \item The system (most likely) has 2 type of RAM:
    \begin{itemize}
      \item A small on-chip RAM (640 KB for i.MX93, 200KB for RK3399)
      \item One or several big (GiB) \textbf{external} memory ((LP)DDR) bank(s)
    \end{itemize}
    \item The on-chip RAM is \textbf{static} and requires little initialization, which can be
          done by the BootROM.
    \item The \textbf{external} RAM is \textbf{dynamic}, it and its controller needs to be initialized
    \item At startup, we must start running from the on-chip RAM, so at least the \textbf{external}
          RAM init logic must have a small memory footprint.
    \item The same goes for using NAND flash as a boot medium.
    \item Modern bootloaders have a lot of features, and therefore are too big.
          Our i.MX93 U-Boot is 850 kiB
    \item We need to split up the boot: a small loader will initialize the \textbf{external} RAM, and
          load the main bootloader there
  \end{itemize}
\end{frame}

\begin{frame}{Secondary Program Loader (SPL)}
  \begin{itemize}
    \item Usually started by the bootROM, also at EL3
    \item In ARM terms, this is \textbf{BL2}
    \item Runs either from SRAM or directly from flash
    \item Initializes DRAM and DRAM controller, usually using firmware
    \item Loads the \textbf{proper} bootloader
    \item Can be stripped down to save space
    \item On TrustZone systems, actually loads TF-A, as well, which returns to U-Boot
    \item Some SoCs (e.g. Rockchip) split DRAM initialization into another
          loader, the TPL.
  \end{itemize}
\end{frame}

\begin{frame}{Boot sequence refresher}
\includegraphics[width=0.7\textwidth]{common/armv8-boot-sequence-common-names.pdf}
\end{frame}

\begin{frame}{The "proper" bootloader}
  \begin{itemize}
    \item On TrustZone platforms, started by TF-A, running at NS-EL2
    \item In ARM terms, this is \textbf{BL32}
    \item Runs from \textbf{external} DRAM
    \item More complex, has support for more peripherals and filesystems
    \item Goal is to load and start the kernel with proper boot arguments
    \item In case of secure boot, verifies the kernel's signature
  \end{itemize}
\end{frame}

\subsection{(Das) U-Boot}

\begin{frame}{\href{https://u-boot.org/}{U-Boot}}
  \begin{itemize}
    \item U-Boot is the de-facto default embedded systems bootloader
    \item Has extensive SoC and boards support
    \item Supports building a \href{https://elixir.bootlin.com/u-boot/v2026.01/A/ident/CONFIG_TPL}{TPL}
          and \href{https://elixir.bootlin.com/u-boot/v2026.01/A/ident/CONFIG_SPL}{SPL}
  \end{itemize}
\end{frame}

\begin{frame}{U-Boot image formats}
  \begin{itemize}
    \item U-Boot needs to load several components:
    \begin{itemize}
      \item the kernel
      \item on ARM, the Device Tree Blob (DTB)
      \item on ARM, SPL needs to load TF-A
      \item on ARM, optionally OP-TEE
      \item optionally, an external ramdisk
    \end{itemize}
    \item We could install all those components separately, and have U-Boot load each
    \item Having one or two images bundling all components would be more practical
    \item This is the role of U-Boot's
          \href{https://elixir.bootlin.com/u-boot/v2026.01/source/tools/mkimage.c}{mkimage}
  \end{itemize}
\end{frame}

\begin{frame}{U-boot image formats}
  \begin{itemize}
    \item U-boot and mkimage support various
          \href{https://elixir.bootlin.com/u-boot/v2026.01/source/boot/image.c\#L141}{image formats}
    \item Some are SoC vendor-specific
    \item U-Boot can only boot 3:
    \begin{itemize}
      \item \href{https://elixir.bootlin.com/u-boot/v2026.01/C/ident/IMAGE\_FORMAT\_LEGACY}{legacy image},
            which is concatenated files with a simple
            \href{https://elixir.bootlin.com/u-boot/v2026.01/source/include/image.h\#L324}{header}
      \item \href{https://elixir.bootlin.com/u-boot/v2026.01/C/ident/IMAGE\_FORMAT\_FIT}{FIT (Flattened Image Tree)}
      \item \href{https://elixir.bootlin.com/u-boot/v2026.01/C/ident/IMAGE\_FORMAT\_ANDROID}{Android boot image},
            which uses this \href{https://source.android.com/docs/core/architecture/bootloader/boot-image-header}{header}
    \end{itemize}
    \item FIT is the current format for non-Android devices
  \end{itemize}
\end{frame}

\begin{frame}{Flattened Image Tree (FIT)}
  \begin{itemize}
    \item The FIT image \href{https://fitspec.osfw.foundation/}{specification} was split off of U-Boot
    \item FIT images are Flattened Device Trees/
          \href{https://devicetree-specification.readthedocs.io/en/stable/flattened-format.html}{Device Tree Blobs}
          that respect additional constraints.
    \item They are essentially containers for sub-\href{https://fitspec.osfw.foundation/\#images-node}{images}
    \item Each of these images can include a
          \href{https://fitspec.osfw.foundation/\#image-signature-nodes}{signature}
    \item The FIT image can also list
          \href{https://fitspec.osfw.foundation/\#configuration-nodes}{configurations},
          which list combinations of the images that can be used.
    \item Configurations can be \href{https://fitspec.osfw.foundation/\#configuration-signature-nodes}{signed}.
          The \code{sign-images} property will list the images to include in the signature.
    \item Support must be enabled via \projconfig{u-boot}{CONFIG_FIT}
  \end{itemize}
\end{frame}

\begin{frame}{FIT generation}
  \begin{itemize}
    \item U-Boot uses \projdir{u-boot}{tools/binman} to generate images
    \item This relies on a "configuration file" in the form of a DTS
    \begin{itemize}
      \item that file is named \code{<pattern>-u-boot.dtsi}
      \item the file should be in \code{arch/<arch>/boot/dts}
      \item as documented in \projfile{u-boot}{tools/binman/binman.rst}, binman will look for the following
            files, in order:
      \begin{itemize}
      \item \code{<dts>-u-boot.dtsi} where <dts> is the base name of the .dts file
      \item \code{<}\projconfig{u-boot}{CONFIG_SYS_SOC}\code{>}\code{-u-boot.dtsi}
      \item \code{<}\projconfig{u-boot}{CONFIG_SYS_CPU}\code{>}\code{-u-boot.dtsi}
      \item \code{<}\projconfig{u-boot}{CONFIG_SYS_VENDOR}\code{>}\code{-u-boot.dtsi}
      \item \code{u-boot.dtsi}
      \end{itemize}
    \end{itemize}
  \end{itemize}
\end{frame}

\begin{frame}[fragile]
  \frametitle{Example: \projfile{u-boot}{arch/arm/dts/imx93-u-boot.dtsi}}
  \begin{minted}[fontsize=\tiny]{c}
/ { binman: binman { multiple-images; }; };
...
&binman {
        u-boot-spl-ddr {
                align = <4>;
                align-size = <4>;
                filename = "u-boot-spl-ddr.bin";
                pad-byte = <0xff>;
                u-boot-spl     { filename = "u-boot-spl.bin"; align-end = <4>; };
                ddr-1d-imem-fw { filename = "lpddr4_imem_1d_v202201.bin"; align-end = <4>; type = "blob-ext"; };
                ...
                ddr-2d-dmem-fw { filename = "lpddr4_dmem_2d_v202201.bin"; align-end = <4>; type = "blob-ext"; };
        };

        spl { filename = "spl.bin";
              mkimage { args = "-n spl/u-boot-spl.cfgout -T imx8image -e 0x2049A000"; blob { filename = "u-boot-spl-ddr.bin"; }; };
        };

        u-boot-container {
                filename = "u-boot-container.bin";
                mkimage { args = "-n u-boot-container.cfgout -T imx8image -e 0x0"; blob { filename = "u-boot.bin"; }; };
        };

        imx-boot {
                filename = "flash.bin";
                pad-byte = <0x00>;
                spl: blob-ext@1 { filename = "spl.bin"; offset = <0x0>; align-size = <0x400>; align = <0x400>; };
                uboot: blob-ext@2 { filename = "u-boot-container.bin"; };
        };
};
  \end{minted}
\end{frame}

\begin{frame}[fragile]
  \frametitle{Configuring FIT verification}
  \begin{itemize}
    \item The feature is documented in \projfile{u-boot}{doc/usage/fit/signature.rst}
    \item \projconfig{u-boot}{CONFIG_FIT_SIGNATURE} must be set
    \begin{itemize}
      \item this will disable \projconfig{u-boot}{CONFIG_LEGACY_IMAGE_FORMAT}, as this format cannot be signed
      \item it is not enough for signatures to be \textbf{required}
    \end{itemize}
    \item The signing key itself will need to be loaded using a \code{.dtsi} file
    \begin{itemize}
      \item This file can be generated using \projfile{u-boot}{tools/key2dtsi.py}
          \begin{minted}{c}
tools/key2dtsi.py --required-image fit_signing_key.pub fit_signing_key.dtsi
            \end{minted}
    \end{itemize}
    \item You can then either
    \begin{itemize}
      \item \code{#include} your \code{fit_signing_key.dtsi} from your main device tree
      \item Use \projconfig{u-boot}{CONFIG_DEVICE_TREE_INCLUDES} to do it for you
    \end{itemize}
  \end{itemize}
\end{frame}

\begin{frame}[fragile]
  \frametitle{fit\_signing\_key.dtsi}
  \begin{minted}{c}
      / {
        signature {
            key-fit_signing_key {
                key-name-hint = "fit_signing_key";
                algo = "sha256,rsa4096";
                rsa,num-bits = <4096>;
                rsa,modulus = [bd 32 f6 a6 5d f7 9a ed \
[...]
                               7c 0b 2f 8e 8f d0 4d 95];
                rsa,exponent = [00 00 00 00 00 01 00 01];
                rsa,r-squared = [bc 5b f8 07 15 a2 36 92 \
[...]
                                 49 1f da e8 b9 74 07 3a];
               rsa,n0-inverse = <0x7bec6a43>;
               required = "image";
            };
        };
    };
  \end{minted}
\end{frame}

\begin{frame}[fragile]
  \frametitle{Configuring FIT signing}
  \begin{itemize}
    \item Now that U-Boot is configured to verify FIT signature we need to:
    \begin{itemize}
      \item pack the kernel and DTB into a FIT
      \item sign the FIT image
    \end{itemize}
    \item Both can be done by \href{https://elixir.bootlin.com/u-boot/v2026.01/source/tools/mkimage.c}{mkimage}
    \begin{minted}{c}
tools/mkimage -k keys_dir -f kernel_fit.its kernel_fit.itb -r
    \end{minted}
    \item The \code{-r} option marks the key as required.
  \end{itemize}
\end{frame}

\begin{frame}[fragile]
  \frametitle{fit.its}
  \begin{minted}[fontsize=\tiny]{c}
/dts-v1/;
/ {
        description = "Image for Linux Kernel";
        images {
                kernel {
                        description = "Linux Kernel";
                        data = /incbin/("Image.gz");
                        ...
                        compression = "gzip";
                        load =  <0xDEADBEEF>;
                        entry = <0xDEADBEEF>;
                        hash-1 { algo = "sha256";};
                };
                fdt {
                        description = "fdt";
                        data = /incbin/("platform.dtb");
                        ...
                };
        };
        configurations {
                default = "config-1";
                config-1 {
                        description = "Linux configuration";
                        kernel = "kernel";
                        fdt = "fdt";
                        signature-1 {
                                algo = \"sha256,rsa2048\";
                                key-name-hint = "fit_signing_key";
                                sign-images = "fdt", "kernel";
                        };

                };
        };
};
  \end{minted}
\end{frame}

\subsection{Last stage: rootFS verification}

\begin{frame}{Root filesystem verification}
  \begin{itemize}
    \item We have verified software up to the kernel, what about userland?
    \item We want to check the integrity of the root filesystem
    \begin{itemize}
      \item The system cannot modify the root filesystem, so it might as well
            be mounted read-only.
      \item We need a hash of the rootFS but a rootFS is usually several orders
            of magnitude larger than the kernel.
    \end{itemize}
    \item The hash will be the linchpin of the verification scheme, so it needs to be signed
    \begin{itemize}
      \item We cannot put the signing key onto the system, so the signature
            must be generated off the system
    \end{itemize}
    \item The kernel must verify the rootFS \textbf{before} it loads anything from it,
          so this will introduce a significant delay in the boot process.
  \end{itemize}
\end{frame}

\begin{frame}{The device mapper (dm)}
  \begin{itemize}
    \item The device mapper will wrap a block device and create a new block device
    \item It defines several types of wrapped devices, called "targets". they include:
    \begin{itemize}
      \item dm-crypt: \kdochtml{admin-guide/device-mapper/dm-crypt}
      \item dm-integrity: \kdochtml{admin-guide/device-mapper/dm-integrity}
      \item dm-verity: \kdochtml{admin-guide/device-mapper/dm-verity}
    \end{itemize}
    among others
    \item These targets can be stacked on top of each other
  \end{itemize}
\end{frame}

\begin{frame}{dm-verity}
  \begin{itemize}
    \item Read-only target of the device-mapper
    \item Needs the \kconfig{CONFIG_DM_VERITY} option enabled
    \item Splits the block device into \textbf{blocks}
    \item Each block then gets hashed to give the leaves of a hash tree
    \item The rest of the tree is built iteratively by aggregating hashes into
          blocks of the data block size, and hashing them.
    \item After \code{math.ceil(math.log(num_blocks, block_size/hash_size))} iterations,
          we are left with one single hash: the \textbf{root hash}
    \item The tree can have maximum \ksym{DM_VERITY_MAX_LEVELS} (63 in 6.18) levels
  \end{itemize}
\end{frame}

\begin{frame}
  \frametitle{dm-verity}
  \center \includegraphics[width=0.8\textwidth]{slides/security-secure-boot-part-2/VerityHashing.pdf}
\end{frame}

\begin{frame}{dm-verity: verification}
  \begin{itemize}
    \item On \code{read()}, dm-verity will perform a block I/O. Before completing it,
          it will call \kfunc{verity_verify_io}
    \item For each block of the I/O, it will
    \begin{itemize}
      \item Lookup whether the block was already verified in the \ksym{validated_blocks} bitfield cache
      \item If not, dm-verity will walk the hash tree down from the root, and for each level, recalculate
            the hashes.
      \item It will then hash the \textbf{actual block data} and compare the hash to the hash from the tree
            that was used in the calculations.
      \item The behaviour on error depends on the device configuration, but can be set to \ksym{DM_VERITY_MODE_PANIC}
    \end{itemize}
  \end{itemize}
\end{frame}

\begin{frame}{dm-verity: setup}
  \begin{itemize}
    \item To properly setup the \code{verity} device, the kernel will need the following parameters
          (described in \kdochtml{admin-guide/device-mapper/verity}):
    \begin{itemize}
      \item \code{version}
      \item \code{dev}
      \item \code{hash_dev}
      \item \code{data_block_size}
      \item \code{hash_block_size}
      \item \code{num_data_blocks}
      \item \code{hash_start_block}
      \item \code{algorithm}
      \item \code{digest}
      \item \code{salt}
    \end{itemize}
    \item This assumes that the root hash has already been calculated, and the hash tree generated.
    \item These arguments can be passed on the kernel command line
    \item The value of the root hash is \textbf{sensitive}, but not \textbf{secret}
    \item Its \textbf{integrity} must be protected, not its \textbf{confidentiality}
  \end{itemize}
\end{frame}

\begin{frame}{Storing the root hash}
  \begin{itemize}
    \item The parameters can be stored in an on-disk verity header
    \begin{itemize}
      % Changing e.g. the data block size will break the device, but not disclose information
      \item Modifying most parameters will be a functional problem, but not a security issue
      % The root hash is the keystone of the whole thing
      \item This is clearly not true for the root hash
    \end{itemize}
    \item The header needs to be signed, or at least the root hash
    \begin{itemize}
      \item The kernel can do that, using \kconfig{CONFIG_DM_VERITY_VERIFY_ROOTHASH_SIG}
      \item It will use the kernel's trusted keyring (see \kdochtml{security/keys/trusted-encrypted})
            via the \kfunc{sys_add_key} system call
      \item The signature will be verified using:
      \begin{itemize}
        \item the trusted keyring built into the kernel,
        \item the secondary keyring if \kconfig{DM_VERITY_VERIFY_ROOTHASH_SIG_SECONDARY_KEYRING}
              is enabled.
        \item the platform keyring if \kconfig{DM_VERITY_VERIFY_ROOTHASH_SIG_PLATFORM_KEYRING}
              is enabled
      \end{itemize}
    \end{itemize}
  \end{itemize}
\end{frame}

\begin{frame}{crypt/veritysetup}
  \begin{itemize}
    \item Interacting with the device can be done using cryptesetup's
          \href{https://gitlab.com/cryptsetup/cryptsetup/-/blob/main/src/veritysetup.c}{veritysetup}
    \begin{itemize}
      \item \code{veritysetup format}
      \item \code{veritysetup open}
      \item \code{veritysetup close}
    \end{itemize}
    \item The root hash signature can be passed using the \code{----root-hash-signature} option
  \end{itemize}
\end{frame}

\begin{frame}{Setup using an initramFS}
  \begin{itemize}
    \item \code{veritysetup} is a userland tool, so it needs a userland to run
    \item We can use an initramFS, see \kdochtml{filesystems/ramfs-rootfs-initramfs}
    \item To avoid breaking our secure boot chain, it must be \textbf{signed}
    \item Fortunately, this is one of the images that can be included in the FIT: \code{FIT_RAMDISK_PROP}
    \item We can then even store the root hash as a file in the initramFS
    \item If we use kernel verification, the signature can be passed from a file
  \end{itemize}
\end{frame}

\begin{frame}{Setup at kernel init}
  \begin{itemize}
    \item The kernel supports boot arguments to create mapped devices: \code{dm-mod.create}
    \item This must be enabled using \kconfig{CONFIG_DM_INIT}
    \item The root hash and salt are then going to be passed as boot arguments
    \item Unfortunately, there is no boot argument to load the signature into the keyring,
          so we can no longer use \kconfig{CONFIG_DM_VERITY_VERIFY_ROOTHASH_SIG}
    \item This means that the kernel command line is now security sensitive, and must be signed
    \begin{itemize}
      \item The FIT spec include a \code{cmdline} property, but it is not implemented in U-Boot
      \item We can use a U-Boot \code{script} to set the \code{bootargs} environment variable.
            The script can be added to the FIT image as a script entry
    \end{itemize}
  \end{itemize}
\end{frame}
